\originalchapter{8.512公开2004年补充作业}
\psource{8.512 Spring 2009作业页另列2004年10套；下表按原卷位置指向已经给出的同题完整题面与解答，再补独有问题。}
\begin{center}\begin{tabular}{lll}\toprule 2004套号&原题&本册定位\\\midrule
1&1、2&习题\ref{p:512-1-1}、\ref{p:512-1-2}\\
2&1、2&习题\ref{p:512-2-1}、\ref{p:512-2-2}（倒格标签改为带标签）\\
3&1、2&习题\ref{p:512-4-1}、\ref{p:512-4-2}\\
4&1&习题\ref{p:512-6-1}\\
6&1&习题\ref{p:512-9-1}\\
7&1&习题\ref{p:512-9-2}\\
8&1、2&习题\ref{p:512-10-1}、\ref{p:512-11-1}\\
9&1--4&习题\ref{p:512-12-1}中对应1--4\\\bottomrule\end{tabular}\end{center}
上述对应保留原套号与年份，不把一次解答重复统计成新增内容。2004套2的$G$标签是折叠带的另一表示，物理条件一致；2004套9与2009套12的个别自旋上下标底稿排版有混乱，均以所列明确分裂能与自洽关系核查。
\originalsection{2004作业4：核磁弛豫（3月11日截止）}
\begin{exercise}[原题2(a)--(b)]
核自旋为两态系统$I$，与电子在原点接触耦合$H_{\rm hf}=A\vct I\cdot\vct S(0)$，$S(r)=\Psi_\alpha^\dagger(r)\vct\sigma_{\alpha\beta}\Psi_\beta(r)$。(a)核自旋初始向上，证明弛豫率可用$\int\avg{S^+(t)S^-(0)}\cos\omega_nt\dd t$表示，原题约定前因子$2A^2/\hbar^2$；$\omega_n$远小于电子尺度可取零。(b)求自由费米气低温率并证明$1/(T_1T)\propto[N(E_F)]^2$，说明超导相干因子的角色。原书阅读指针不收录。
\end{exercise}
\begin{solution}
核翻转耦合来自$A(I^-S^++I^+S^-)/2$，Golden rule对电子热初末态求和，转为时间相关的Fourier变换；纵向回到平衡的$T_1$还合并上翻与下翻率。因原题自旋与$S^\pm$归一不同于常见$S=\sigma/2$，固定其前因子后使用所给定义，不能无条件复制另一教材的数值2。

自由电子$S^+(0)$在同一位置将一自旋态变成另一，谱权重为$\int\rho_\uparrow(E)\rho_\downarrow(E+\hbar\omega_n)f(E)[1-f(E+\hbar\omega_n)]\dd E$。小$\omega_n$及平滑态密度下$\int f(1-f)\dd E=\kb T$，故$1/T_1=\mathcal C(A^2/\hbar)\rho_\uparrow(E_F)\rho_\downarrow(E_F)\kb T$，$\mathcal C$由原自旋归一固定，Korringa温度与态密度结论不依该记号。超导态须将电子算符作Bogoliubov展开，核自旋通道的$uu'+vv'$型相干因子与能隙态密度共同决定低温行为；原题只要求阅读商业书了解改动，未给独立需证明的完整超导率题面。
\end{solution}
\originalsection{2004作业5：节点超导与响应（3月18日截止）}
\begin{exercise}[原题1(a)--(d)]
二维$d$波隙$\Delta(\phi)=\Delta_0\cos2\phi$，正常费米面圆形，原题取$\hbar=\kb=1$，能谱为
\begin{equation}
E(k)=\sqrt{v_F^2(|k|-k_F)^2+\Delta(\phi)^2}.
\end{equation}
(a)找4节点、线性化$E=\sqrt{v_F^2k_1^2+v_2^2k_2^2}$、求$v_2$及每节点每自旋态密度。(b)证明$\rho_s(T)/m=\rho_s(0)/m-(2\log2/\pi)(v_F/v_2)T$。(c)在极端II类单涡旋外$\xi_0\ll R\ll\lambda$，Doppler位移$v_F\cdot[\nabla\theta/2+eA/c]$使零能态密度非零，分别求沿$x$与$(1,1)$接近涡核的结果。(d)三角涡晶格中估计$c_V=\alpha\sqrt H\,T$及$\alpha$量级；原题所列实验论文只作参考链接。
\end{exercise}
\begin{solution}
(a)节点$\phi=\pi/4+n\pi/2$，近节点$\Delta\simeq\pm2\Delta_0\delta\phi=\pm(2\Delta_0/k_F)k_2$，故$v_2=2\Delta_0/k_F$。缩放$p_1=v_Fk_1,p_2=v_2k_2$，$\int\delta(E-|p|)\dd^2p/(2\pi)^2v_Fv_2=E/(2\pi v_Fv_2)$。

(b)热准粒子的反向电流给$\delta(\rho_s/m)=-2\sum_{\rm nodes}\int v_x^2[-f'(E)]\dd^2k/(2\pi)^2$，节点的$x$分量平方$v_F^2/2$；四节点并双自旋归一后系数为$-(2/\pi)(v_F/v_2)\int_0^\infty E[-f'(E)]\dd E$。积分分部给$\int f(E)\dd E=T\log2$，得到线性律。

(c)流动动量$\vct p_s\simeq\hat\phi/(2R)$，每节点的$D(0)=|v_F\hat n\cdot p_s|/(2\pi v_Fv_2)$，自旋另乘2。$R\parallel x$时四节点投影绝对值皆$1/\sqrt2$，总$D(0,R)=\sqrt2/(\pi v_2R)$；$R\parallel(1,1)$时两节点投影零、两节点为1，总$1/(\pi v_2R)$。角依赖是节点结构，不是标量正常态密度的简单平移。

(d)涡间距$r_H\sim\sqrt{\Phi_0/H}$，把$1/R$在一个胞内积分给平均$\overline D(0)\sim1/(v_2r_H)\propto\sqrt H$，因此费米型低温$C\propto\overline D(0)T$。用圆胞近似$R_c=\sqrt{\Phi_0/(\pi H)}$，角平均$\overline{\sum_n|\sin(\phi_n-\phi)|}=8/\pi$，得$\overline D\simeq8/(\pi^2v_2R_c)$，所以按该态密度与常用热归一$\alpha\simeq8\sqrt\pi\kb^2/(3v_2\sqrt{\Phi_0})$。真实三角格胞、核区和速度单位恢复改变数值常数，$\sqrt H\,T$是本题稳健的低场低温量级，不宣称这个粗略常数精确。
\end{solution}
\begin{exercise}[原题2]
列完美金属、无序金属、洁净超导、无序超导的实部横纵$K$，分别取静态$\omega=0,q\to0$与动态$q=0,\omega\to0$，共16个量，并给主要非零贡献。
\end{exercise}
\begin{solution}
采用SI核$j=-KA$，$D=ne^2/m,D_s(T)=n_s(T)e^2/m$，相位已作规范一致的纵向调整。表中洁净体系取单带Galilean不变连续模型、电流与总动量成正比且没有动量弛豫；无序体系取有限弹性散射时间$\tau$。温度低于$T_c$，先固定$T$再取表中静态或动态极限。16个实部的极限与先导项如下；零不表示整条频率函数恒零。
\begin{center}\begin{tabular}{lllll}\toprule 系统&静态$K_T$&静态$K_L$&动态$K_T$&动态$K_L$\\\midrule
完美金属&$-\chi_Dq^2/\mu_0$&0&$D$&$D$\\
无序金属&$-\chi_Dq^2/\mu_0$&0&$D(\omega\tau)^2$&$D(\omega\tau)^2$\\
洁净超导&$D_s(T)+O(q^2)$&0&$D+O(\omega^2)$&$D+O(\omega^2)$\\
无序超导&$D_s(T)+O(q^2)$&0&$D_s(T)+O(\omega^2)$&$D_s(T)+O(\omega^2)$\\\bottomrule\end{tabular}\end{center}
无序正常金属全复核$K=-\ii\omega D\tau/(1-\ii\omega\tau)$，实部$D(\omega\tau)^2/[1+(\omega\tau)^2]$，虚部首阶$-\omega\sigma_0$，故动态实极限为0。洁净Galilean体系因$\vct j\propto\vct P$且$\vct P$守恒，$q=0$动态核是$D$，即使$0<T<T_c$仍然如此；静态Meissner核却是$D_s(T)<D$。两极限不可互换。有限弛豫的无序超导在$\omega\to0$只保留凝聚体刚度$D_s(T)$，正常准粒子的正则电导给$-\ii\omega\sigma_{\rm reg}$并在实部极限消失；$T=0$时洁净Galilean体系$D_s(0)=D$。静态纵$A$是纯规范，若固定序参量相位而不让它响应，会得到错误非零值。
\end{solution}
\originalsection{2004作业7：易轴磁振子（4月13日截止）}
\begin{exercise}[原题2(a)--(c)，公开文件引Ashcroft--Mermin第723页题5]
各向异性铁磁$H=-\frac12\sum_{RR'}[J_z(R-R')S_z(R)S_z(R')+J_\perp(R-R')\vct S_\perp(R)\cdot\vct S_\perp(R')]$，$J_z>J_\perp>0$。(a)证明全极化与单自旋波仍为本征态，并给相对各向同性的能量抬升。(b)证明低温磁化偏离饱和按激活指数小。(c)解释二维无连续自旋序的红外论证为何失效。原PDF的抬升式两项都排为$J_z$，应由哈密顿量独立判断。
\end{exercise}
\begin{solution}
全上态横向交换翻转不能作用，故本征。单翻转可在格点间传播，纵向代价$S\sum_RJ_z(R)$，横向迁移给$-SJ_\perp(k)$，所以$E_k=S[J_z(0)-J_\perp(k)]$（此处$J(0)$表示Fourier零波矢）。相对各向同性$SJ_\perp(0)-SJ_\perp(k)$，抬升为$\Delta=S\sum_R[J_z(R)-J_\perp(R)]>0$，不是原底稿误排的零。

低温$E_k\simeq\Delta+Ak^2$，磁振子数$\int\e^{-E_k/T}\dd^dk\propto T^{d/2}\e^{-\Delta/T}$，得到激活磁化减小。二维连续对称模型$\int\dd^2k/k^2$红外发散；易轴打开隙且剩离散翻转对称，积分$\int\dd^2k/(\Delta+Ak^2)$红外有限，所以那条禁止论证不再成立，但仅“不再禁止”不是单凭此积分就证明某温度必有序。
\end{solution}
\originalsection{2004作业10：共振杂质的总电荷（5月6日截止）}
\begin{exercise}[原题1--6]
非相互作用$d$轨道$\epsilon_d$与传导带通过$V_k$杂化，$\Sigma(E)=\sum_k|V_k|^2/(E-\epsilon_k+\ii0)$，$G_d=[E-\epsilon_d-\Sigma]^{-1}$，$\Gamma=-\Sigma''$，相移由$\tan\delta=\Gamma/(\epsilon_d+\Sigma'-E)$连续选支。(1)用$\operatorname{Im}\log G_d$表达相移。(2)积分$\partial_E\log G_d$写为$G_d$和$\Sigma'$的式。(3)求$\delta(E_F)$与$n_d$的关系。(4)以$G_k-G_k^0$求传导电荷改变。(5)证明Friedel规则含总电荷。(6)估计窄共振$\Gamma\ll E_F$时忽略$\Sigma$能量依赖的误差。公开底稿个别式的负号与$1/\pi$归一需按Green函数定义重核。
\end{exercise}
\begin{solution}
(1)推迟$G_d$的相位在$[-\pi,0]$，使空共振在低能的相移为0可取$\delta=\pi+\operatorname{Im}\log G_d$，连续追踪支，不直接用$\arctan$主值跳变。(2)$\partial_E\log G_d=-G_d(1-\partial_E\Sigma)$，积分后
\begin{equation}
\delta(E_F)=\pi n_d+\operatorname{Im}\int_{-\infty}^{E_F}G_d(E)\partial_E\Sigma(E)\dd E,\qquad n_d=-\frac1\pi\operatorname{Im}\int^{E_F}G_d\dd E.
\end{equation}
这个符号从微分恒等式直接得出，不能沿用排版相反号却同时声称匹配总电荷。

(4)Dyson方程$G_k-G_k^0=|V_k|^2(G_k^0)^2G_d$，而$\partial_E\Sigma=-\sum_k|V_k|^2(G_k^0)^2$，故$\sum_k(G_k-G_k^0)=-G_d\partial_E\Sigma$。因此$\delta n_c=-\pi^{-1}\operatorname{Im}\int\sum_k(G_k-G_k^0)\dd E=\pi^{-1}\operatorname{Im}\int G_d\partial_E\Sigma\dd E$，(5)$\delta(E_F)=\pi(n_d+\delta n_c)$。

(6)若杂化带函数在能量$E_F$或带宽$D$上变化，$\partial_E\Sigma=O(\Gamma/D)$，共振窗口的积分修正相对为该小尺度。常$\Gamma,\Sigma'$的宽带近似使$\delta n_c=0$，才得到$\delta/\pi=n_d$。接近带边、伪隙或强能量依赖时窄峰本身不保证修正小，必须检查实际$\Sigma'$尺度。
\end{solution}
